\bookchapter{Глава 8: Связь: SMS, звонки и соцсети}

\chapterintro{Связь — самая первая и главная задача телефона. Сегодня она превратилась из простых звонков и SMS в сложную экосистему с мессенджерами, соцсетями и электронной почтой.}

\booksection{Продвинутые функции приложения «Телефон»}

Современное приложение «Телефон» давно переросло простой номеронабиратель: в нём есть запись разговоров, блокировка спама, умное определение номеров и другие полезные инструменты.

\booksubsection{Запись разговоров (обратите внимание на законность):}

\begin{itemize}
  \item Зачем: фиксировать важные разговоры — для рабочих протоколов, интервью, обещаний клиентов или сохранения доказательств при спорах.
  \item Юридическая подсказка: в Китае запись разговора как доказательство допустима при согласии собеседника или для защиты своих законных прав.
  \item Односторонняя запись против взаимного согласия: с юридической точки зрения запись, сделанная «без согласия собеседника», не считается автоматически недействительной. Согласно разъяснениям Верховного народного суда, если способ получения записи не нарушает серьёзно права других лиц, не противоречит закону и нормам морали, её можно предъявить как доказательство. Однако в ряде провинций суды строже относятся к «тайной записи» в гражданских делах; в уголовных делах требования к законности обычно выше.
  \item Региональные различия: суды Пекина, Шанхая, Гуандуна и ряда других мест лояльнее относятся к односторонней записи как к гражданскому доказательству. Но записи, затрагивающие коммерческую тайну или личную жизнь, требуют особой осторожности в любой провинции — их могут исключить как нарушение права на неприкосновенность частной жизни.
  \item Как записать: на Android в экране звонка обычно есть явная кнопка «Запись»; iOS не предоставляет встроенную запись разговоров — нужны сторонние приложения или Bluetooth-гарнитура с функцией записи.
\end{itemize}

\begin{tcbwarn}{Красные линии при использовании записей}
\begin{itemize}
  \item Нельзя использовать записи для шантажа, распространения или публичного оскорбления других людей.
  \item Материалы, затрагивающие частную жизнь третьих лиц (например, случайно записанный чужой разговор), нельзя свободно распространять.
  \item В важных случаях рекомендуется в начале разговора устно сообщить: «Этот разговор будет записан» — это и законно, и усиливает доказательную силу.
\end{itemize}

\end{tcbwarn}

\booksubsection{Углублённо о борьбе со спам-звонками:}

\begin{itemize}
  \item Принцип: приложение «Телефон» обращается к обширной базе номеров по сети и при входящем звонке показывает «спам-звонок», «мошеннический звонок» или название компании.
  \item Как включить: на Android откройте Телефон > Настройки > Блокировка спама; на iOS — Настройки > Телефон > Блокировка и идентификация. Подробные пути переключений по брендам и рекомендуемые опции — в таблице брендов главы 13 «Практические системные инструменты».
  \item Основные платформы маркировки номеров: в Китае данные о маркировке номеров в основном предоставляют три сервиса — Tencent Mobile Manager, 360 Mobile Guard и Dianhuabang (Tapir). Операторы (China Mobile / China Unicom / China Telecom) при показе «имени владельца» и в функции блокировки спама тоже часто опираются на эти базы. Базы маркировки производителей телефонов (Huawei, Xiaomi, OPPO, vivo) по большей части работают в партнёрстве с теми же сервисами.
  \item Что делать при ошибочной маркировке: если ваш обычный рабочий номер ошибочно пометили как «доставка еды», «курьерская служба» или «спам-звонки», подайте апелляцию через следующие каналы: Dianhuabang: зайдите на платформу апелляций по номерам (haoma.tk.cn) и отправьте документы, подтверждающие личность.
  \item 360: подайте заявку через функцию «Апелляция по номеру» в приложении 360 Mobile Guard.
  \item Tencent: подайте заявку на странице «Апелляция по номеру» в мини-приложении «Tencent Mobile Manager».
  \item Три крупных оператора: абоненты China Mobile могут отправить SMS с текстом «KTFSR» на номер 10086, чтобы включить защиту от спам-звонков; абонентам China Telecom стоит подписаться на официальный аккаунт «Tianyi Anti-Harassment»; абоненты China Unicom настраивают защиту в приложении «Unicom Phone Manager».
\end{itemize}

Продвинутая защита: включение у оператора «блокировки частых вызовов» (например, ограничение высокочастотных вызовов China Mobile) помогает автоматически блокировать голосовые обзвонные звонки ИИ.

Типичные возможности встроенной блокировки: Идентификация / умная блокировка: автоматическая маркировка и блокировка спама, мошенничества и рекламных звонков по базе номеров.

Чёрный / белый список: вручную добавляйте известные спам-номера в чёрный список; номера из контактов, курьеров и доставки — в белый, чтобы избежать ложных срабатываний.

Фильтрация SMS по ключевым словам: сообщения со словами «вы выиграли», «кредит», «перейдите по ссылке» и т. п. автоматически попадают в папку блокировки.

iOS: «Неизвестные вызовы беззвучно»: звонки не из контактов беззвучно уходят в голосовую почту и не мешают. Удобно при большом количестве спама, но следите, чтобы не пропустить важный звонок с незнакомого номера.

Связка с «Не беспокоить»: ночью или на совещании разрешайте звонки только из контактов, включайте по расписанию — качество сна сразу улучшится.

Как сочетать полезные программы: сначала используйте встроенную в систему блокировку, затем добавьте Национальный центр борьбы с мошенничеством и антиспам-сервис оператора; для очистки подойдёт лёгкий инструмент вроде Qingzhuo, но приложений каждого типа (безопасность/очистка) достаточно одного — подробнее в главе 13.

\booksubsection{Что такое VoLTE/VoNR и как их включить:}

\begin{itemize}
  \item VoLTE (Voice over LTE): технология голосовых звонков поверх сети 4G. Во время звонка телефон не переключается на 2G/3G, можно одновременно разговаривать и пользоваться интернетом, звук чище, а соединение устанавливается быстрее.
  \item VoNR (Voice over NR): голосовые звонки поверх 5G — следующий этап после VoLTE. Задержка ниже, качество звука лучше, поддерживается видеосвязь высокой чёткости. Нужны 5G-телефон, 5G-тариф и покрытие 5G в вашем районе.
  \item Как включить: на iOS — Настройки > Сотовая связь > Параметры сотовых данных > Голос и данные, включите «VoLTE»; на Android — Настройки > Мобильная сеть > Мобильные данные > Дополнительно, найдите переключатель «VoLTE-высокое качество». Операторы обычно включают его по умолчанию; если во время звонка интернет пропадает — скорее всего, VoLTE не активирован.
\end{itemize}

\begin{tcbinfo}{Полезно знать}
Если во время звонка интернет внезапно тормозит или пропадает, скорее всего VoLTE/VoNR выключен и телефон упал на 2G/3G. После включения можно одновременно разговаривать и смотреть видео.

\end{tcbinfo}

\booksection{SMS и iMessage/RCS}

Обычные SMS не устарели: они остаются надёжным каналом для кодов подтверждения, банковских уведомлений и официальных сообщений.

\begin{itemize}
  \item iMessage: сообщения между устройствами Apple через интернет, отображаются голубыми пузырями. Поддерживают качественные фото и видео, уведомления о прочтении, полностью бесплатны и имеют сквозное шифрование.
  \item RCS: протокол следующего поколения, который продвигают мировые операторы и Google. По функциям похож на iMessage или WeChat, но работает на основе телефонного номера.
  \item Совместимость iMessage и RCS: начиная с iOS 18 Apple поддерживает стандарт RCS — между iPhone и Android можно отправлять фото высокого качества, уведомления о прочтении и статус набора текста. Прощай, разрыв «зелёных пузырей». Но нужно, чтобы RCS был включён у обоих операторов, а версии систем были достаточно новыми.
\end{itemize}

\booksubsection{Развёртывание RCS в Китае:}

\begin{itemize}
  \item China Mobile: раньше других внедрила RCS, в некоторых провинциях поддерживается сервис «5G-сообщения» с мультимедийными сообщениями и интерактивными карточками. Нужен поддерживающий телефон и подключение в салоне связи.
  \item China Telecom: услуга предоставляется через «China Telecom 5G Messages»; в части тарифов она включена по умолчанию, а давним пользователям нужно отправить специальную SMS-команду для подключения.
  \item China Unicom: разворачивает «5G-сообщения» в крупных городах, но поддержка устройств неравномерна — во многих Android-моделях переключатель «Улучшенные сообщения» или «5G-сообщения» приходится включать вручную в приложении SMS.
  \item Практика: в Китае RCS распространяется медленнее ожидаемого из-за сильной экосистемы WeChat и неоднородной поддержки устройств. В повседневном общении по-прежнему доминирует WeChat, а RCS чаще играет роль «официальных рассылок» (банковские уведомления, государственные сообщения).
\end{itemize}

\begin{tcbwarn}{Безопасность кодов подтверждения}
\begin{itemize}
  \item Не пересылайте коды подтверждения: любой, кто просит «назвать код» или «переслать SMS», может быть мошенником — включая «службу поддержки», «курьера» и «полицию». Код подтверждения = доступ к вашим деньгам.
  \item Фишинговые SMS: остерегайтесь ссылок в сообщениях вида «ETC просрочен», «медкарта заблокирована», «компенсация за утерянную посылку». Официальные организации не запрашивают номер карты и пароль по ссылке из SMS.
  \item Идентификация номера: официальные банковские SMS обычно приходят с коротких номеров, начинающихся на 95, или с 1069 + подпись компании. «Банковское уведомление» с обычного 11-значного мобильного номера почти наверняка мошенничество.
  \item Резерв второго фактора: не меняйте без необходимости номер, привязанный к важным аккаунтам. Если номер менять нужно, сначала отвяжите все ключевые аккаунты — иначе новый владелец сможет сбросить пароли по SMS-коду.
\end{itemize}

\end{tcbwarn}

\booksection{WeChat/QQ: управление национальными мессенджерами}

WeChat и QQ давно переросли простые мессенджеры и стали «суперприложениями» для общения, платежей и получения информации.

\booksubsection{Резервное копирование и перенос переписки:}

\begin{enumerate}
  \item Войдите в WeChat на компьютере (PC или Mac).
  \item Нажмите меню в левом нижнем углу > [Перенос и резервное копирование] > [Резервное копирование и восстановление].
  \item Выберите «Сохранить переписку на компьютер» и на телефоне отметьте нужные чаты.
  \item Перенос на новый телефон: в WeChat на старом телефоне [Я] > [Настройки] > [Основные] > [Перенос и резервное копирование переписки] > [Перенос] > [Перенести на телефон/планшет], отсканируйте QR-код — готово.
\end{enumerate}

\booksubsection{Углублённая очистка места:}

\begin{itemize}
  \item Обычный вход: WeChat [Я] > [Настройки] > [Основные] > [Место хранения] — можно очистить кэш и точечно удалить переписку по дате и типу файлов.
  \item Очистка внутри чата: откройте чат > «…» в правом верхнем углу > [Переписка] > [Управление местом] — можно удалить фото, видео и файлы конкретного чата, сохранив текст.
  \item Очистка оригиналов: самые тяжёлые — оригинальные фото и видео из групповых чатов. В [Управление местом] отметьте «Оригиналы» и удалите пакетом.
  \item Отключите автозагрузку: [Настройки] > [Основные] > [Фото, видео, файлы и звонки], выключите «Автозагрузку», чтобы в Wi-Fi незаметно не кэшировались групповые файлы.
  \item Кэш мини-приложений: в [Обзор] > [Мини-приложения] удерживайте неиспользуемые и удалите — освободится локальный кэш.
  \item Сжатое резервное копирование: при резервном копировании с компьютера отметьте «Сжать» — несколько гигабайт переписки можно уменьшить до 30–50\% от исходного размера.
\end{itemize}

\booksubsection{Приёмы управления группами:}

\begin{itemize}
  \item Объявление группы: владелец/админ публикует в «…» > [Объявление группы] — все получают принудительное уведомление. Подходит для важных сообщений; появляется автоматически при входе новых участников.
  \item Задачи группы: владелец может закрепить сообщение как «Задачу группы» — у всех участников появится красная отметка, после выполнения она исчезает. Удобно для регистрации на события и сбора задач.
  \item Цепочки (jielong): введите в поле ввода «\#jielong» («цепочка»), чтобы запустить цепочку — удобно для коллективного заказа еды и сбора записей. Результаты автоматически сводятся вверху группового чата и не засоряют экран.
  \item Групповой эфир: «+» в правом верхнем углу чата > [Групповой эфир] — трансляция только для этой группы. Подходит для корпоративного обучения, семейных встреч, закрытого тестирования продукта.
  \item Свернуть группу: в списке чатов смахните имя группы влево > [Свернуть] — ненужные группы уйдут в «Свёрнутые группы», уведомления только при @упоминании вас.
  \item Только владелец/админ может @всех: в группах от 500 человек обычные участники не могут упомянуть всех — это снижает шум.
\end{itemize}

\booksubsection{Тонкая настройка прав в «Моментах»:}

\begin{itemize}
  \item Видимость по группам: при публикации нажмите «Кто может видеть», выберите «Только некоторые» или «Скрыть от кого-то» — на основе заранее созданных меток («Коллеги», «Семья», «Клиенты»).
  \item Объём видимости для друзей: [Я] > [Настройки] > [Права друзей] > [Моменты] > [Объём видимости], можно выбрать «Последние 3 дня / месяц / полгода / всё». Рекомендуется «полгода» — баланс приватности и общения.
  \item Скрыть чужую запись: удерживайте запись в «Моментах» > [Не видеть его/её моменты] или [Отписаться] — автор не узнает.
  \item Управление рекламой: [Я] > [Настройки] > [О WeChat] > [Личные данные и права] > [Управление персонализированной рекламой] — можно отключить поведенческие рекомендации.
  \item Десять записей для незнакомцев: [Права друзей] > [Моменты] > [Разрешить незнакомцам видеть 10 записей]. Отключите, чтобы чужие не листали вашу историю.
\end{itemize}

\booksection{Сравнение видеозвонков}

У разных инструментов видеосвязи разные сильные стороны — выбор под сценарий делает общение удобнее.

\begin{itemize}
  \item WeChat Video: общенациональное покрытие, подходит всем возрастам. Для повседневного общения с близкими и быстрых созвонов. До 20 человек (нужно мини-приложение «Групповое видео»). Минус — при плохой сети качество заметно падает.
  \item FaceTime: экосистема Apple, отличное качество и сквозное шифрование. С iOS 15 можно приглашать Android/Web-пользователей по ссылке, но ведущий должен быть на устройстве Apple. Удобно для семейного общения между разными устройствами.
  \item Zoom: выбор для международных встреч, бесплатно до 100 участников 40 минут. Экран, виртуальный фон, групповые комнаты — удобно для трансграничной и удалённой работы.
  \item Tencent Meeting: выбор для офиса в Китае, тесная интеграция с WeChat, вход по QR. До 300 участников (платная версия), облачная запись, субтитры и протоколы — для корпоративного обучения и больших совещаний.
  \item Рекомендации: для личного общения 1-на-1 — WeChat; между пользователями Apple — FaceTime; для деловых встреч — Zoom или Tencent Meeting (в Китае — Tencent, за рубежом — Zoom).
\end{itemize}

\booksubsection{Требования к каналу и расход трафика:}

\begin{itemize}
  \item WeChat Video: около 500 Кбит/с–1 Мбит/с; стандартное качество \textasciitilde{}3–5 МБ в минуту, HD \textasciitilde{}8–12 МБ.
  \item FaceTime: SD \textasciitilde{}1 Мбит/с, 720p \textasciitilde{}2 Мбит/с, 1080p от 4 Мбит/с. Расход \textasciitilde{}5–15 МБ в минуту.
  \item Zoom/Tencent Meeting: 720p \textasciitilde{}1,5–3 Мбит/с, 1080p 3–5 Мбит/с. Расход \textasciitilde{}300–800 МБ в час; у ведущего многопользовательской встречи трафика больше.
  \item Совет: для длительных видеозвонков используйте Wi-Fi; в мобильной сети можно сознательно снизить качество (отключите «HD-видео» в настройках приложения), чтобы сэкономить трафик.
\end{itemize}

\booksubsection{Советы по улучшению качества видео:}

\begin{itemize}
  \item Свет: сидите спиной к окну, лицом к источнику света (лучше всего естественный), избегайте контрового света, из-за которого лицо темнеет.
  \item Угол: камера чуть выше уровня глаз — так меньше виден второй подбородок; в горизонтальном режиме держите телефон двумя руками.
  \item Звук: наушники убирают эхо и свист; в тихом месте не включайте громкую связь.
  \item Сеть: предпочтительнее Wi-Fi; в 4G/5G стойте ближе к окну, чтобы не терять сигнал; в многопользовательской встрече можно скрыть видео всех, кроме говорящего.
  \item Размытие фона: Tencent Meeting и Zoom поддерживают ИИ-размытие, чтобы спрятать беспорядок. Носите одежду спокойных цветов, не сливающуюся с фоном, иначе возможен «эффект зелёного экрана».
\end{itemize}

\begin{tcbinfo}{Совет}
При лагах сети лучше отключить видео и оставить голос, чем сразу сбрасывать. Перед важной встречей проверьте канал функцией «Проверка сети» в Tencent Meeting.

\end{tcbinfo}

\booksection{Электронная почта: настройка и использование на телефоне}

\begin{itemize}
  \item Добавление аккаунта: iOS — Настройки > Почта > Учётные записи > Добавить учётную запись; на Android обычно есть приложение «Электронная почта».
  \item Приёмы: задайте разные способы уведомлений для разных почтовых ящиков; используйте свайпы для быстрого архивирования, удаления и отметки писем.
\end{itemize}

\booksubsection{Чем отличаются IMAP, POP3 и Exchange:}

\begin{itemize}
  \item POP3: письма скачиваются с сервера и по умолчанию удаляются, остаются только локально. Подходит для одного устройства и офлайн-чтения. Минус — на другом устройстве не будет истории; сейчас не рекомендуется.
  \item IMAP: письма всегда хранятся на сервере, локально только кэш. Устройства (телефон/компьютер/планшет) синхронизированы — прочитано в одном месте, везде прочитано. Основной рекомендуемый протокол.
  \item Exchange (включая ActiveSync): протокол Microsoft, кроме почты синхронизирует календарь, контакты и задачи. Лучший выбор для корпоративной почты. Поддерживается Outlook, NetEase и Tencent Enterprise Mail.
  \item Рекомендация: для личной почты — IMAP; для рабочей — сначала Exchange, затем IMAP; POP3 по возможности не используйте, если нет особых задач архивации.
\end{itemize}

\booksubsection{Настройка корпоративной почты:}

\begin{enumerate}
  \item iOS: [Настройки] > [Почта] > [Учётные записи] > [Добавить учётную запись] > выберите «Microsoft Exchange» или «Другие».
  \item Android: в приложении «Электронная почта» или Outlook выберите «Добавить аккаунт», введите полный адрес и пароль.
  \item Заполнение сервера (пример Exchange): адрес обычно mail.домен-компании.com или outlook.office365.com, домен — короткое имя компании.
  \item Если появляется «Не удаётся проверить сервер», введите вручную: входящий сервер (IMAP) порт 993, шифрование SSL; исходящий (SMTP) порт 465 или 587, шифрование TLS.
  \item После входа включите в [Настройках аккаунта] «Push» для получения писем в реальном времени и отключите «Предпросмотр письма», чтобы защитить коммерческую тайну.
\end{enumerate}

\begin{tcbinfo}{Подсказка}
Корпоративная почта обычно требует от администратора «пароль клиента» или «код авторизации IMAP/SMTP» — пароль от веб-почты не подойдёт. При ошибке входа сначала сгенерируйте код в веб-версии почты, затем введите его на телефоне.

\end{tcbinfo}

\begin{tcbpractice}{Практика: сохранить переписку WeChat на компьютер}
\begin{enumerate}
  \item Запустите резервное копирование на компьютере: откройте WeChat на ПК > меню в левом нижнем углу > [Перенос и резервное копирование] > [Резервное копирование и восстановление].
  \item Выберите сохранение на компьютер: нажмите «Сохранить переписку на компьютер».
  \item На телефоне выберите содержимое для резервной копии и дождитесь завершения.
\end{enumerate}

\end{tcbpractice}
